\documentclass[main.tex]{subfiles}

\begin{document}\setWPP
\lsection{Limity}
\begin{multicols}{2}
\qaw{l'Hospitalovo pravidlo}{
Slouží k výpočtu limit neurčitých výrazů typu $\left[\frac{0}{0}\right]$ nebo $\left[\frac{\infty}{\infty}\right]$.
Nechť $f(x)$ a $g(x)$ mají v okolí bodu $a$ derivace a existuje limita podílu jejich derivací:
$$\lim_{x \to a} \frac{f(x)}{g(x)} \stackrel{l'H}{=} \lim_{x \to a} \frac{f'(x)}{g'(x)}$$
\textbl{Důležité upozornění:} Derivujeme čitatele i jmenovatele zvlášť (nejedná se o derivaci podílu $\left(\frac{f}{g}\right)'$).
}
\columnbreak
\qaw{Důležité tabulkové limity}{
Základní vztahy pro výpočty bez l'Hospitalova pravidla ($t \to 0$):
\begin{align*}
    \lim_{t \to 0} \frac{\sin t}{t} &= 1, \qquad \lim_{t \to 0} \frac{e^t - 1}{t} = 1 \\
    \lim_{t \to 0} \frac{\ln(1 + t)}{t} &= 1, \qquad \lim_{t \to \infty} \left(1 + \frac{1}{t}\right)^t = e
\end{align*}
Tyto limity se často aplikují spolu se substitucí $t = g(x)$.
}
\end{multicols}
\wbreak
\begin{multicols}{2}
\qaw{Hornerovo schéma -- postup}{
Hornerovo schéma slouží k dělení polynomu $P(x) = a_n x^n + \dots + a_1 x + a_0$ kořenovým činitelem $(x - c)$.
\textbl{Kroky výpočtu:}
\begin{enumerate}[label={\color{WPP}\textbf{\arabic*.}}]
    \item Do prvního řádku zapíšeme koeficienty $a_n, a_{n-1}, \dots, a_0$.
    \item Do prvního sloupce druhého řádku dáme číslo $c$ (předpokládaný kořen).
    \item První koeficient $a_n$ opíšeme dolů.
    \item Další čísla počítáme podle pravidla: \\
    $\text{nová hodnota} = c \cdot (\text{předchozí výsledek}) + \text{koeficient nad ním}$.
    \item Poslední číslo je zbytek po dělení (je-li $c$ kořen, zbytek $= 0$).
\end{enumerate}
}
\columnbreak
\qaw{Názorný příklad Hornerova schématu}{
Dělíme $x^3 - 5x^2 + 3x + 9$ číslem $c = 3$:
$$\begin{array}{c|c c c c}
    & 1 & -5 & 3 & 9 \\
\hline
x = 3 & 1 & -2 & -3 & \mathbf{0}
\end{array}$$
\textbl{Postup:}
\begin{itemize}
    \item Opíšeme $1$.
    \item $3 \cdot 1 + (-5) = -2$.
    \item $3 \cdot (-2) + 3 = -3$.
    \item $3 \cdot (-3) + 9 = \mathbf{0}$ (zbytek je $0$).
\end{itemize}
Výsledek reprezentuje polynom $x^2 - 2x - 3$.
}
\end{multicols}
\wbreak
\begin{multicols}{2}
\qaw{Metoda 1: l'Hospitalovo pravidlo}{
Máme zadánu limitu:
$$\lim_{x \to 1} \frac{e^x - e}{x - 1}$$
Po dosazení $x=1$ dostáváme typ $\left[\frac{0}{0}\right]$. Použijeme l'Hospitalovo pravidlo (derivujeme čitatele i jmenovatele):
\begin{align*}
    \lim_{x \to 1} \frac{e^x - e}{x - 1} &\stackrel{l'H}{=} \lim_{x \to 1} \frac{(e^x - e)'}{(x - 1)'} \\
    &= \lim_{x \to 1} \frac{e^x - 0}{1 - 0} \\
    &= \lim_{x \to 1} e^x = e^1 = e
\end{align*}
}
\columnbreak
\qaw{Metoda 2: Úprava na základní limitu}{
Vytkneme z čitatele konstantu $e$:
$$\lim_{x \to 1} \frac{e(e^{x-1} - 1)}{x - 1} = e \cdot \lim_{x \to 1} \frac{e^{x-1} - 1}{x - 1}$$
Substituce: $t = x - 1$. Pro $x \to 1$ jde $t \to 0$:
$$e \cdot \lim_{t \to 0} \frac{e^t - 1}{t}$$
Využijeme známou tabulkovou limitu $\lim_{t \to 0} \frac{e^t - 1}{t} = 1$:
$$e \cdot 1 = e$$
}
\end{multicols}
\begin{center}
\begin{tikzpicture}
\begin{axis}[
    xmin=-0.5, xmax=2.5,
    ymin=0, ymax=6,
    axis lines=middle,
    axis line style={color=antiDP, thick},
    tick style={color=antiDP},
    tick label style={color=antiDP, font=\tiny},
    xlabel={\color{antiDP}$x$}, ylabel={\color{antiDP}$y$},
    width=0.42\textwidth, 
    height=5cm,
    unit vector ratio=1.5 1, % Adjusts the ratio between x and y unit lengths
    samples=100
]
    \addplot[domain=-0.5:0.99, color=WPP, thick]{(exp(x)-exp(1))/(x-1)};
    \addplot[domain=1.01:2.5, color=WPP, thick]{(exp(x)-exp(1))/(x-1)};
    \addplot[only marks, mark=o, mark size=2pt, color=WPP, thick] coordinates {(1, 2.718)};
\end{axis}
\end{tikzpicture}
\end{center}
\newpage
\begin{multicols}{2}
\qaw{Metoda 1: l'Hospitalovo pravidlo}{
Máme zadánu limitu:
$$\lim_{x \to 1} \frac{1 - \sqrt{x}}{\ln x}$$
Po dosazení $x=1$ dostáváme neurčitý výraz typu $\left[\frac{0}{0}\right]$. Použijeme l'Hospitalovo pravidlo (derivujeme čitatele i jmenovatele):
\begin{align*}
    \lim_{x \to 1} \frac{1 - \sqrt{x}}{\ln x} &\stackrel{l'H}{=} \lim_{x \to 1} \frac{\left(1 - x^{\frac{1}{2}}\right)'}{(\ln x)'} \\
    &= \lim_{x \to 1} \frac{-\frac{1}{2\sqrt{x}}}{\frac{1}{x}} \\
    &= \lim_{x \to 1} \left( -\frac{x}{2\sqrt{x}} \right) \\
    &= -\frac{1}{2\cdot 1} = -\frac{1}{2}
\end{align*}
}
\columnbreak
\qaw{Metoda 2: Substituce a tabulková limita}{
Upravíme výraz v čitateli i jmenovateli:
$$\lim_{x \to 1} \frac{1 - \sqrt{x}}{\ln x} = \lim_{x \to 1} \frac{1 - \sqrt{x}}{\ln(1 + (x - 1))}$$
Substituce: $t = x - 1 \implies x = 1 + t$. Pro $x \to 1$ jde $t \to 0$:
$$\lim_{t \to 0} \frac{1 - \sqrt{1+t}}{\ln(1 + t)}$$
Rozšíříme čitatele i jmenovatele výrazem $(1 + \sqrt{1+t})$:
\begin{align*}
    &= \lim_{t \to 0} \frac{1 - (1+t)}{\ln(1 + t) \cdot (1 + \sqrt{1+t})} \\
    &= \lim_{t \to 0} \left( \frac{-t}{\ln(1 + t)} \cdot \frac{1}{1 + \sqrt{1+t}} \right)
\end{align*}
Využijeme známou limitu $\lim_{t \to 0} \frac{\ln(1+t)}{t} = 1$:
$$-1 \cdot \frac{1}{1 + \sqrt{1}} = -\frac{1}{2}$$
}
\end{multicols}
\begin{center}
\begin{tikzpicture}
\begin{axis}[
    xmin=0, xmax=5,
    ymin=-1.8, ymax=0.5,
    axis lines=middle,
    axis line style={color=antiDP, thick},
    tick style={color=antiDP},
    tick label style={color=antiDP, font=\tiny},
    xlabel={\color{antiDP}$x$}, ylabel={\color{antiDP}$y$},
    width=0.42\textwidth, height=4.5cm,
    unit vector ratio=1 1,
    samples=150
]
    \addplot[domain=0.05:0.99, color=WPP, thick]{(1-sqrt(x))/ln(x)};
    \addplot[domain=1.01:5, color=WPP, thick]{(1-sqrt(x))/ln(x)};
    \addplot[only marks, mark=o, mark size=2pt, color=WPP, thick] coordinates {(1, -0.5)};
\end{axis}
\end{tikzpicture}
\end{center}
\begin{multicols}{2}
\qaw{Detailní postup (vytknutí $x$)}{
Počítáme limitu v záporném nekonečnu:
$$\lim_{x \to -\infty} \frac{\sqrt{2x^2 - 1}}{x - 2}$$
Z čitatele pod odmocninou vytkneme $x^2$:
$$\lim_{x \to -\infty} \frac{\sqrt{x^2 \left(2 - \frac{1}{x^2}\right)}}{x - 2}$$
\textbl{Pozor na znaménko:} Protože $x < 0$, platí $\sqrt{x^2} = |x| = -x$:
$$\lim_{x \to -\infty} \frac{-x \cdot \sqrt{2 - \frac{1}{x^2}}}{x \left(1 - \frac{2}{x}\right)}$$
Po vykrácení $x$ v čitateli a jmenovateli:
$$\lim_{x \to -\infty} \frac{-\sqrt{2 - \frac{1}{x^2}}}{1 - \frac{2}{x}}$$
Jelikož $\frac{1}{x^2} \to 0$ a $\frac{2}{x} \to 0$ pro $x \to -\infty$:
$$\frac{-\sqrt{2 - 0}}{1 - 0} = \texthl{-\sqrt{2}}$$
}
\columnbreak
\qaw{Stručné shrnutí a chyták}{
\textbl{Častá chyba:}
Pokud by ses překlikl a použil $\sqrt{x^2} = x$ (což platí jen pro $x > 0$), vyšlo by ti nesprávně $+\sqrt{2}$.
\textbl{Porovnání limit pro $\pm\infty$:}
\begin{align*}
    \lim_{x \to +\infty} \frac{\sqrt{2x^2 - 1}}{x - 2} &= \texthl{+\sqrt{2}} \\
    \lim_{x \to -\infty} \frac{\sqrt{2x^2 - 1}}{x - 2} &= \texthl{-\sqrt{2}}
\end{align*}
\textbl{Pravidlo pro asymptoty:}
Graf funkce $f(x) = \frac{\sqrt{2x^2 - 1}}{x - 2}$ má tedy dvě různé horizontální asymptoty v závislosti na směru do nekonečna.
}
\end{multicols}
\begin{center}
\begin{tikzpicture}
\begin{axis}[
    xmin=-8, xmax=8,
    ymin=-3, ymax=3,
    axis lines=middle,
    axis line style={color=antiDP, thick},
    tick style={color=antiDP},
    tick label style={color=antiDP, font=\tiny},
    xlabel={\color{antiDP}$x$}, ylabel={\color{antiDP}$y$},
    width=0.42\textwidth, height=5cm,
    unit vector ratio=2 1,
    samples=200
]
    % Záporná větev (x <= -1/sqrt(2))
    \addplot[domain=-8:-0.71, color=WPP, thick]{sqrt(2*x^2 - 1)/(x - 2)};
    
    % Kladné větve (x >= 1/sqrt(2), vyhýbáme se vertikální asymptotě v x=2)
    \addplot[domain=0.71:1.9, color=WPP, thick]{sqrt(2*x^2 - 1)/(x - 2)};
    \addplot[domain=2.1:8, color=WPP, thick]{sqrt(2*x^2 - 1)/(x - 2)};

    % Horizontální asymptoty y = -sqrt(2) a y = +sqrt(2)
    \addplot[domain=-8:8, color=antiDP, dashed]{-1.414};
    \addplot[domain=-8:8, color=antiDP, dashed]{1.414};
\end{axis}
\end{tikzpicture}
\end{center}
\newpage
\begin{multicols}{2}
\qaw{Rozklad polynomů (Hornerovo schéma)}{
Počítáme limitu:
$$\lim_{x \to 3} \frac{x^3 - 5x^2 + 3x + 9}{x^2 - 6x + 9}$$
Po dosazení $x=3$ dostáváme typ $\left[\frac{0}{0}\right]$. Rozložíme čitatele pomocí Hornerova schématu pro kořen $x=3$:
$$\begin{array}{c|c c c c}
    & 1 & -5 & 3 & 9 \\
\hline
x = 3 & 1 & -2 & -3 & 0
\end{array}$$
Čitatel: $(x - 3)(x^2 - 2x - 3) = (x - 3)^2 (x + 1)$
Jmenovatel: $x^2 - 6x + 9 = (x - 3)^2$
}
\columnbreak
\qaw{Výpočet limity}{
Dosadíme rozložené tvary do limity a zkrátíme kořenový činitel $(x - 3)^2$:
\begin{align*}
    \lim_{x \to 3} \frac{x^3 - 5x^2 + 3x + 9}{x^2 - 6x + 9} &= \lim_{x \to 3} \frac{(x - 3)^2 (x + 1)}{(x - 3)^2} \\
    &= \lim_{x \to 3} (x + 1) \\
    &= 3 + 1 = \texthl{4}
\end{align*}
\textbl{Výsledek:} Limita je rovna \texthl{4}.
}
\end{multicols}
\begin{center}
\begin{tikzpicture}
\begin{axis}[
    xmin=-2, xmax=5,
    ymin=-1, ymax=6,
    axis lines=middle,
    axis line style={color=antiDP, thick},
    tick style={color=antiDP},
    tick label style={color=antiDP, font=\tiny},
    xlabel={\color{antiDP}$x$}, ylabel={\color{antiDP}$y$},
    width=0.42\textwidth, height=4.5cm,
    unit vector ratio=1 1,
    samples=100
]
    \addplot[domain=-2:2.95, color=WPP, thick]{(x^3 - 5*x^2 + 3*x + 9)/(x^2 - 6*x + 9)};
    \addplot[domain=3.05:5, color=WPP, thick]{(x^3 - 5*x^2 + 3*x + 9)/(x^2 - 6*x + 9)};
    \addplot[only marks, mark=o, mark size=2pt, color=WPP, thick] coordinates {(3, 4)};
\end{axis}
\end{tikzpicture}
\end{center}
\wpage
\end{document}